\documentclass[11pt]{article}
\usepackage[margin=1in]{geometry}
\usepackage{amsmath,amssymb,amsthm}
\usepackage{algorithm}
\usepackage{algpseudocode}
\usepackage{booktabs}
\usepackage{hyperref}

\newtheorem{definition}{Definition}
\newtheorem{theorem}{Theorem}
\newtheorem{lemma}[theorem]{Lemma}
\newtheorem{property}{Property}

\title{\textbf{Mathematical Specification, Robbins-Monro Convergence Proof, and L2 Regularization Dynamics of Stochastic Gradient Descent (ALGO-NN-31)}}
\author{Team Scaibu \\ \small Email: \texttt{jaydeep.9664920749@gmail.com} \\ \small Architecture and Core Engine Team}
\date{\today}

\begin{document}
\maketitle

\begin{abstract}
Stochastic Gradient Descent (SGD) updates parameters iteratively along negative stochastic subgradient trajectories evaluated on mini-batches of empirical training data. We formalize the Robbins-Monro stochastic approximation framework, prove convergence to stationary points on $L$-smooth non-convex objectives, provide line-by-line pseudocode matching the reference implementation, and derive complexity.
\end{abstract}

\section{Notation and Mathematical Preliminaries}
\label{sec:notation}

Let $\theta \in \mathbb{R}^P$ denote the parameter vector for $P$ neural network weights, and $\mathcal{L}(\theta) = \mathbb{E}_{\xi}[f(\theta; \xi)]$ the population risk with stochastic mini-batch gradient $\mathbf{g}_t = \nabla f(\theta_t; \xi_t)$.

\subsection{SGD Parameter Update Mechanics}
\paragraph{Intuitive Meaning:} Instead of computing gradients across the entire million-sample dataset (which is slow and memory-intensive), SGD takes a fast step using a small random mini-batch. The noise in the mini-batch gradient helps the model escape bad local minima.

\paragraph{Formal Mathematical Definition:}
For learning rate $\eta > 0$ and weight decay coefficient $\lambda \ge 0$:
\begin{equation}
\label{eq:sgd_update}
\mathbf{g}_{\mathrm{eff}, t} = \mathbf{g}_t + \lambda \theta_t, \quad \theta_{t+1} = \theta_t - \eta \mathbf{g}_{\mathrm{eff}, t}.
\end{equation}

\paragraph{Worked Numerical Example:}
Let $\theta = [1.0, -2.0]$, gradient $\mathbf{g} = [0.1, 0.4]$, $\eta = 0.01, \lambda = 0.05$:
$\mathbf{g}_{\mathrm{eff}} = [0.1 + 0.05(1.0), 0.4 + 0.05(-2.0)] = [0.15, 0.30]$.
$\Delta\theta = 0.01 \times [0.15, 0.30] = [0.0015, 0.0030]$.
$\theta_{\mathrm{next}} = [1.0 - 0.0015, -2.0 - 0.0030] = [0.9985, -2.0030]$.

\subsection{Summary of Symbols}
\begin{itemize}
    \item $P \in \mathbb{N}$: Number of trainable parameter coordinates.
    \item $\boldsymbol{\theta} \in \mathbb{R}^P$: Current parameter vector.
    \item $\mathbf{g} \in \mathbb{R}^P$: Stochastic mini-batch gradient vector.
    \item $\eta \in \mathbb{R}_{>0}$: Learning rate step size ($\mathrm{lr}$).
    \item $\lambda \in \mathbb{R}_{\ge 0}$: L2 weight decay regularization coefficient.
    \item $\boldsymbol{\theta}_{\mathrm{next}} \in \mathbb{R}^P$: Updated parameter vector.
    \item $\|\Delta\boldsymbol{\theta}\|_2 \in \mathbb{R}_{\ge 0}$: Parameter displacement step norm.
\end{itemize}

\section{Problem Definition and Core Concepts}
\label{sec:problem_definition}

\subsection{Robbins-Monro Convergence Conditions}
Stationary convergence $\lim_{t \to \infty} \mathbb{E}[\|\nabla f(\theta_t)\|_2^2] = 0$ requires the step size sequence $\{\eta_t\}$ to satisfy:
\begin{equation}
\label{eq:robbins_monro}
\sum_{t=1}^\infty \eta_t = \infty \quad \text{(infinite travel distance)}, \quad \sum_{t=1}^\infty \eta_t^2 < \infty \quad \text{(variance suppression)}.
\end{equation}

\section{Algorithmic Specification}
\label{sec:algorithm}

Algorithm~\ref{alg:sgd} specifies the evaluation routine matching \texttt{impl.py}.

\begin{algorithm}[H]
\caption{Stochastic Gradient Descent (SGD) Parameter Update}
\label{alg:sgd}
\begin{algorithmic}[1]
\Require Parameters $\boldsymbol{\theta} \in \mathbb{R}^P$, gradients $\mathbf{g} \in \mathbb{R}^P$, learning rate $\eta > 0$, weight decay $\lambda \ge 0$.
\Ensure Updated parameters $\boldsymbol{\theta}_{\mathrm{next}} \in \mathbb{R}^P$, step norm $S_{\mathrm{norm}}$.
\State Initialize $\boldsymbol{\theta}_{\mathrm{next}} \gets [], \quad step\_sq \gets 0.0$
\For{$i = 1$ to $P$}
    \State $g_{\mathrm{eff}} \gets g_i + \lambda \cdot \theta_i$
    \State $\Delta_i \gets \eta \cdot g_{\mathrm{eff}}$
    \State $\theta_{\mathrm{new}} \gets \theta_i - \Delta_i$
    \State Append $\theta_{\mathrm{new}}$ to $\boldsymbol{\theta}_{\mathrm{next}}$
    \State $step\_sq \gets step\_sq + \Delta_i^2$
\EndFor
\State $S_{\mathrm{norm}} \gets \sqrt{step\_sq}$
\State \Return $\boldsymbol{\theta}_{\mathrm{next}}, \quad S_{\mathrm{norm}}$
\end{algorithmic}
\end{algorithm}

\section{Theoretical Guarantees and Convergence Bounds}
\label{sec:guarantees}

\begin{theorem}[Non-Convex Stationary Convergence of SGD]
\label{thm:sgd_convergence}
\textbf{Assumptions:} Let $f(\theta)$ be $L$-smooth ($f(y) \le f(x) + \nabla f(x)^\top (y-x) + \frac{L}{2}\|y-x\|_2^2$), lower-bounded by $f^* > -\infty$, with bounded gradient variance $\mathbb{E}[\|\mathbf{g}_t - \nabla f(\theta_t)\|_2^2] \le \sigma^2$ and learning rate $\eta_t = \frac{\eta_0}{\sqrt{T}}$.
\textbf{Guarantee:} The average squared gradient norm converges at rate $\min_{t \le T} \mathbb{E}[\|\nabla f(\theta_t)\|_2^2] \le \mathcal{O}\left(\frac{1}{\sqrt{T}}\right)$.
\textbf{Proof:}
Applying $L$-smoothness to the update step $\theta_{t+1} = \theta_t - \eta \mathbf{g}_t$:
\begin{equation}
f(\theta_{t+1}) \le f(\theta_t) - \eta \nabla f(\theta_t)^\top \mathbf{g}_t + \frac{L\eta^2}{2} \|\mathbf{g}_t\|_2^2.
\end{equation}
Taking conditional expectations $\mathbb{E}[\mathbf{g}_t \mid \theta_t] = \nabla f(\theta_t)$ and $\mathbb{E}[\|\mathbf{g}_t\|_2^2] \le \|\nabla f(\theta_t)\|_2^2 + \sigma^2$:
\begin{equation}
\mathbb{E}[f(\theta_{t+1}) \mid \theta_t] \le f(\theta_t) - \eta \|\nabla f(\theta_t)\|_2^2 + \frac{L\eta^2}{2} \left(\|\nabla f(\theta_t)\|_2^2 + \sigma^2\right) = f(\theta_t) - \eta\left(1 - \frac{L\eta}{2}\right)\|\nabla f(\theta_t)\|_2^2 + \frac{L\eta^2\sigma^2}{2}.
\end{equation}
Choosing $\eta \le \frac{1}{L}$ such that $1 - \frac{L\eta}{2} \ge \frac{1}{2}$:
\begin{equation}
\frac{\eta}{2} \|\nabla f(\theta_t)\|_2^2 \le f(\theta_t) - \mathbb{E}[f(\theta_{t+1})] + \frac{L\eta^2\sigma^2}{2}.
\end{equation}
Summing from $t=1$ to $T$ and dividing by $T \eta$:
\begin{equation}
\frac{1}{T}\sum_{t=1}^T \mathbb{E}[\|\nabla f(\theta_t)\|_2^2] \le \frac{2(f(\theta_1) - f^*)}{T\eta} + L\eta\sigma^2.
\end{equation}
Setting $\eta = \frac{1}{\sqrt{T}}$ yields $\mathcal{O}\left(\frac{1}{\sqrt{T}}\right)$ convergence.
\textbf{Limitations:} In ill-conditioned ravines, SGD oscillates orthogonally to the optimal trajectory, motivating momentum acceleration.
\end{theorem}

\section{Complexity Analysis}
\label{sec:complexity}

\begin{itemize}
    \item \textbf{Time Complexity:} $\mathcal{O}(P)$ sequential elementwise multiplications and subtractions.
    \item \textbf{Space Complexity:} $\mathcal{O}(P)$ for updated parameter storage. Zero optimizer memory states required.
\end{itemize}

\section{Worked Numerical Example}
\label{sec:worked_example}

Let $\theta = [2.0], \mathbf{g} = [0.5], \eta = 0.1, \lambda = 0.0$:
\begin{enumerate}
    \item $g_{\mathrm{eff}} = 0.5 + 0.0 = 0.5$.
    \item $\Delta = 0.1 \times 0.5 = 0.05$.
    \item $\theta_{\mathrm{next}} = 2.0 - 0.05 = 1.95$.
    \item Step norm $= \sqrt{0.05^2} = 0.05$.
\end{enumerate}

\section{Numerical Considerations and Edge Cases}
\label{sec:numerical}

\begin{itemize}
    \item \textbf{Decoupled vs. L2 Weight Decay:} In standard SGD, adding $\lambda \theta$ directly to the gradient is mathematically identical to weight decay $\theta \gets \theta(1 - \eta \lambda)$ because the step multiplier is unscaled by adaptive moments.
\end{itemize}

\begin{thebibliography}{9}
\bibitem{robbins1951} H.~Robbins and S.~Monro, ``A Stochastic Approximation Method,'' \emph{Annals of Mathematical Statistics}, vol.~22, no.~3, pp.~400--407, 1951.
\end{thebibliography}

\end{document}
